%% ===============================================================
\section{Experimental Setup}
\label{sec:experiments}

\subsection{Dataset}

We evaluate on the \textbf{Pitt Ads Dataset}~\cite{hussain2017automatic}, a collection of 3,477 video advertisements. Of these, 1,704 remained available for download (many original YouTube links had been removed). The full pipeline was run on all accessible videos; the benchmark comparison against baselines uses a subset of 484 videos due to API budget constraints.

The accessible videos span diverse characteristics:
\begin{itemize}
    \item \textbf{Duration}: Short ($\leq$20\,s, $n{=}98$), Medium (20--45\,s, $n{=}487$), Long (45--75\,s, $n{=}320$), Extended ($>$75\,s, $n{=}227$)
    \item \textbf{Categories}: Food \& Beverage (298), Lifestyle \& Home (233), Social Issues \& PSA (156), Entertainment \& Media (142), Technology (72), Automotive (64), and others---across 38 fine-grained topic categories
    \item \textbf{Cut Frequency}: Low ($<$0.17 cuts/s, $n{=}562$), Medium (0.17--0.47 cuts/s, $n{=}580$), High ($>$0.47 cuts/s, $n{=}562$), based on tertile binning (range: 0.015--1.066 cuts/s, mean = 0.38)
\end{itemize}

Ground-truth annotations use a standardized 38-category topic taxonomy from the original Pitt Ads annotation protocol.

\subsection{Baselines}

We compare against 9 frame selection methods, all using Gemini-3-Flash for extraction:
\begin{itemize}
    \item \textbf{Uniform-1FPS}: 1 frame per second (industry standard).
    \item \textbf{Random}: Random frame sampling at the same count as Uniform-1FPS.
    \item \textbf{Histogram}: Histogram-based change detection at 100\,ms intervals.
    \item \textbf{ORB}: ORB feature-based change detection.
    \item \textbf{Optical Flow}: Dense optical flow change detection.
    \item \textbf{CLIP Only}: CLIP embedding clustering without hash/LPIPS pre-filtering.
    \item \textbf{K-Means}: K-Means clustering on CLIP embeddings (the clustering component inside our cascade, run standalone without the cascade or TA-NMS).
    \item \textbf{PySceneDetect}: Content-aware scene boundary detection via adaptive thresholding.
    \item \textbf{DSNet}: A deep reinforcement learning approach for diverse video summarization.
\end{itemize}

%% Data based on experiments run on 500-video dataset (Mar 2026)
% PySceneDetect & 20.6 & 0.237 & 71.8 & 76.8 & 4.12 & 0.211 \\
% DSNet         & 200.5 & 2.306 & 65.0 & 69.6 & 3.71 & 0.155 \\
% Gemini Native & --- & --- & 82.4 & 87.2 & 4.11 & --- \\

\subsection{Implementation Details}

Implemented in Python~3.11 with PyTorch~2.1. CLIP ViT-B/32 provides embeddings, Whisper-large-v3 handles transcription, Tesseract~5.3 performs OCR. All experiments use \textbf{Gemini-3-Flash} as the VLM provider. Experiments run on an NVIDIA RTX 4090.

%%PLACEHOLDER: Add Claude/GPT-4V results paragraph once multi-VLM experiment completes.
%%PLACEHOLDER: Add Gemini native video comparison once experiment completes.

\subsection{Metrics}

\textbf{Compression Ratio}: $N / |S|$ where $N$ = total frames at native FPS, $|S|$ = selected frames.
\textbf{Topic Accuracy}: Exact match of predicted topic ID against Pitt Ads ground truth (38 categories).
\textbf{Super-category Accuracy}: Match at grouped super-category level (11 groups).
\textbf{Cost}: USD per video based on Gemini-3-Flash pricing.
\textbf{Pairwise Agreement}: Fraction of videos where two methods produce identical predictions.
